\section*{Abstract}\label{sec:abstract}

Organisations consist of multiple internal systems containing large amounts of data, such as project
tools, documentation, communication channels and information about employees.
This spread of information leads to employees spending unnecessary amounts of time searching for
internal information on a daily basis.

This study therefore investigates how data from internal systems can be represented to provide a
large language model (LLM) with the context it needs to navigate and retrieve internal information,
thereby improving information retrieval within an organisation.
The primary focus is on whether a graph-based data representation in the form of a knowledge graph
gives LLMs better conditions for combining information from multiple sources compared to an
API-based solution.

Two solutions are implemented and evaluated. The graph-based solution stores data from GitHub,
Slack and Confluence in a graph database, while the API-based solution retrieves data directly from
the systems’ APIs. Both solutions are exposed via an MCP server and tested using 15 questions
distributed across three complexity levels.

The results show that the graph-based solution performs better on complex questions that require
linking data from different sources. It achieved correct answers on 80\% of such questions compared
to 20\% for the API-based solution, and averaged 2.4 tool calls compared to 6.
For simpler questions, both solutions performed on average equally well.

The conclusion is that a graph-based representation gives the LLM clear advantages when combining
information from multiple sources, while the API-based solution is sufficient for answering
questions that concern a single source.

\textbf{Keywords:} knowledge graph, graph database, LLM, information retrieval, MCP,
API-based solution
